% ============================================================
%  CopilotScope — Tematy prac dyplomowych
%  Dziesięć tematów prac inżynierskich i magisterskich
%  wyprowadzonych z dziesięciu algorytmów oceny jakości sesji
%  opisanych w quality_measurement_framework.tex, docs/ANALYSIS.md §8
%  i w tabeli "Algorithm status" na stronie projektu.
%
%  Kompilacja: pdflatex thesis_topics.tex  (dwukrotnie)
% ============================================================
\documentclass[11pt,a4paper]{article}

\usepackage[T1]{fontenc}
\usepackage[utf8]{inputenc}
\usepackage[polish]{babel}
\usepackage{lmodern}
\usepackage[margin=1.9cm,top=1.7cm,bottom=1.6cm,headsep=0.5cm]{geometry}
\usepackage{booktabs}
\usepackage{array}
\usepackage{xcolor}
\usepackage{enumitem}
\usepackage{framed}
\usepackage{fancyhdr}
\usepackage{microtype}
\usepackage{parskip}
\usepackage{hyperref}
\usepackage{url}

% \kod{...} — ścieżki/identyfikatory czcionką maszynową, łamane po / _ . - :
% (bez dzielenia wyrazów), żeby długie ścieżki nie wychodziły poza margines.
\urlstyle{tt}
\def\UrlBreaks{\do\/\do\_\do\.\do\-\do\:\do\=\do\,}
\newcommand{\kod}{\begingroup\urlstyle{tt}\Url}

\definecolor{csblue}{HTML}{3C3489}
\definecolor{csrule}{HTML}{5F5E5A}
\definecolor{shadecolor}{HTML}{EEEDFE}
\definecolor{csgrey}{HTML}{4A4A48}

\hypersetup{
  colorlinks=true,
  linkcolor=csblue,
  urlcolor=csblue,
  pdftitle={CopilotScope — tematy prac dyplomowych},
  pdfauthor={CopilotScope Project},
}

\setlist[itemize]{leftmargin=1.1em,topsep=1pt,itemsep=0.5pt,parsep=0pt}
\setlength{\parskip}{3pt}
\setlength{\emergencystretch}{3em}   % ostatnia deska ratunku dla wierszy ze ścieżkami

\pagestyle{fancy}
\fancyhf{}
\renewcommand{\headrulewidth}{0.4pt}
\fancyhead[L]{\footnotesize\color{csgrey}CopilotScope — tematy prac dyplomowych}
\fancyhead[R]{\footnotesize\color{csgrey}\thepage}

% ---- makra układu tematu ----

% \temat{numer}{tytuł}{poziom}{metryczka}
\newcommand{\temat}[4]{%
  \clearpage
  \noindent{\footnotesize\color{csblue}\textsc{Temat #1}\hfill\textsc{#3}}\\[-0.5em]
  \noindent\textcolor{csblue}{\rule{\linewidth}{1.2pt}}\\[0.1em]
  {\large\bfseries #2\par}
  \vspace{0.15em}
  \noindent{\footnotesize\color{csgrey}#4\par}
  \vspace{0.1em}
  \noindent\textcolor{csrule}{\rule{\linewidth}{0.4pt}}
  \vspace{-0.2em}
}

% \blok{nagłówek} — śródtytuł biegnący w linii
\newcommand{\blok}[1]{\vspace{0.2em}\noindent{\bfseries\color{csblue}#1}\par\nopagebreak}

% streszczenie w ramce
\newenvironment{streszczenie}%
  {\vspace{-0.7em}\begin{shaded}\noindent\textbf{Streszczenie.\ }}%
  {\end{shaded}\vspace{-0.5em}}

\title{%
  \textbf{CopilotScope --- tematy prac dyplomowych}\\[0.4em]
  \large Dziesięć propozycji prac inżynierskich i magisterskich\\
  z obszaru pomiaru jakości sesji asystenta programistycznego
}
\author{Projekt CopilotScope}
\date{Sierpień 2026}

\begin{document}
\maketitle
\thispagestyle{fancy}

\vspace{-1em}

\section*{Skąd wzięło się tych dziesięć tematów}
\addcontentsline{toc}{section}{Skąd wzięło się tych dziesięć tematów}

CopilotScope zbiera telemetrię OpenTelemetry z asystentów programistycznych
(GitHub Copilot w~VS~Code, Copilot CLI, Claude Code) i wylicza z~niej złożoną
ocenę jakości sesji w~skali 0--100. Przegląd \emph{dziesięciu} podejść do oceny
jakości rozmowy --- ten sam zestaw, który strona projektu pokazuje w~tabeli
\textit{Algorithm status}, a~\texttt{docs/ANALYSIS.md} §8 rozwija w~formie
przeglądu literatury --- stanowi mapę tego obszaru. Formalne definicje
zaimplementowanych algorytmów (wzory, dowody własności, przykłady liczbowe)
znajdują się w~artykule
\textit{CopilotScope Quality Measurement Framework}
(\texttt{research/articles/quality\_measurement\_framework.tex}), a~ich
odtworzenia w~Pythonie --- w~notatnikach \texttt{research/0X\_*.ipynb}.

Każdy z~dziesięciu algorytmów zamieniono tutaj na \textbf{samodzielny temat pracy
dyplomowej}: jedna strona A4, z~tytułem, streszczeniem, celem, zakresem prac,
metodyką, mierzalnym kryterium akceptacji i~punktem wejścia w~kodzie. Numeracja
jest zgodna z~tabelą \textit{Evaluation algorithms --- implementation status}
w~\texttt{README.md} oraz z~§8a w~\texttt{docs/ANALYSIS.md}, żeby odwołania
między dokumentami się nie rozjeżdżały.

\textbf{Sześć z~dziesięciu algorytmów ma już działającą implementację}
w~katalogu \texttt{Collector/Quality/} --- tam praca dyplomowa polega na
\emph{walidacji i~kalibracji} tego, co dziś działa na ręcznie dobranych
parametrach. \textbf{Cztery nie mają implementacji} --- tam praca polega na
\emph{projekcie i~prototypie}. Jedno i~drugie jest pełnoprawnym tematem; drugie
jest tylko bardziej efektowne, a~pierwsze zwykle bardziej potrzebne.

\begin{table}[h!]
\centering
\small
\caption{Mapa tematów: algorytm $\rightarrow$ temat $\rightarrow$ poziom pracy.}
\begin{tabular}{@{}clll@{}}
\toprule
\# & Algorytm & Stan implementacji & Proponowany poziom \\
\midrule
1  & LLM-as-a-Judge (G-Eval)          & nie zaimplementowany   & magisterska \\
2  & SPUR --- wyuczone rubryki SAT/DSAT & nie zaimplementowany  & magisterska \\
3  & Metryki komponentowe RAG (RAGAS) & nie zaimplementowany   & inżynierska \\
4  & Edit Survival Analysis           & w pełni                & inżynierska \\
5  & Acceptance-weighted throughput   & w pełni                & inżynierska \\
6  & TFRA --- tarcie i naprawy per tura & w pełni              & magisterska \\
7  & Latency-utility model            & w pełni                & inżynierska \\
8  & Ekonomia tokenów i cache         & w pełni                & inżynierska \\
9  & Klasyfikacja frustracji          & uproszczony (report-only) & magisterska \\
10 & Task-completion detection        & częściowe hooki        & inżynierska \\
\bottomrule
\end{tabular}
\end{table}

\blok{Co jest wspólne dla wszystkich dziesięciu}
\begin{itemize}
  \item \textbf{Dane.} \texttt{tools/CopilotScope.Seeder -{}- demo} generuje sesje
        trzech person (\emph{clean}, \emph{error-prone}, \emph{frustrated}) oraz
        zestaw długich, scenariuszowych rozmów (\textit{Redis rate limiter},
        \textit{RDS Postgres upgrade} i~inne). To punkt startowy każdej pracy;
        sesje z~realnego użycia repozytorium (dogfooding) są cenniejsze, ale
        trzeba je najpierw zebrać.
  \item \textbf{Punkt wejścia w~kodzie.} Interfejs \texttt{IInsightAnalyzer}
        w~\texttt{Quality/Insights.cs}: nowy algorytm to jedna klasa plus jedna
        rejestracja w~kontenerze DI --- zero pracy w~interfejsie użytkownika,
        bo panel \textit{Insights} renderuje raporty generycznie.
  \item \textbf{Trzy zasady projektowe, których praca nie powinna łamać po
        cichu} (artykuł, §Introduction): \emph{telemetry-only} (bez sędziego LLM
        w~momencie scoringu), \emph{wyjaśnialność z~konstrukcji} (każdy wynik
        rozkłada się na nazwane składowe z~powodami), \emph{świadomość pewności}
        (brak danych to \texttt{confidence}~=~0, nie fałszywa neutralna ocena).
        Temat, który tę zasadę narusza --- jak \#1 czy \#2 --- musi to nazwać
        wprost i~uzasadnić, a~nie przemycić.
  \item \textbf{Granica etyczna.} CopilotScope świadomie nie przechowuje
        tożsamości osoby przy sesji. Żadna z~tych prac nie ma prowadzić do oceny
        \emph{ludzi} --- przedmiotem pomiaru jest jakość sesji z~asystentem, nie
        wydajność programisty (patrz sekcja \textit{How not to use CopilotScope}
        w~\texttt{README.md}).
\end{itemize}

\blok{Uwaga o poziomie pracy}
Podział na inżynierską i~magisterską jest propozycją, nie wyrocznią. Kryterium
przyjęte tutaj: \emph{magisterska} tam, gdzie sednem jest zaprojektowanie
eksperymentu z~udziałem ludzi (anotacja, miary zgodności) albo wyuczenie modelu
i~wykazanie przewagi nad baseline'em; \emph{inżynierska} tam, gdzie sednem jest
zbudowanie działającego narzędzia i~zmierzenie nim czegoś konkretnego. Każdy
temat inżynierski da się podnieść do magisterskiego, rozszerzając część badawczą
--- w~opisach zaznaczono, którędy.

% ============================================================
\temat{1}{Ocena jakości sesji asystenta programistycznego przez model-sędziego:
projekt rubryki G-Eval i~jej walidacja względem ocen ludzkich}
{Praca magisterska}
{Algorytm \#1 --- LLM-as-a-Judge (G-Eval) \quad·\quad stan: nie zaimplementowany \quad·\quad wymaga agenta-sędziego (Azure AI Foundry)}

\begin{streszczenie}
Model-sędzia ocenia transkrypt sesji względem jawnej rubryki --- poprawność,
kompletność, styl komunikacji --- i~zwraca wynik, który ma być
\emph{porównywalny} ze składanym wskaźnikiem 0--100 wyliczanym dziś wyłącznie
z~telemetrii. Praca nie odpowiada na pytanie „czy G-Eval działa'' (to technika
ugruntowana), tylko na trudniejsze: czy da się go wpiąć w~system, którego
zasadą projektową jest \emph{telemetry-only}, bez zerwania tej zasady --- i~po
jakim koszcie tokenowym przestaje się to opłacać.
\end{streszczenie}

\blok{Kontekst i motywacja}
Sześć algorytmów w~CopilotScope działa na samych metadanych: nie czyta treści
promptów, nie wywołuje modelu, nie wychodzi do sieci. To świadomy wybór, ale ma
cenę --- żaden z~nich nie ocenia, czy odpowiedź asystenta była
\emph{merytorycznie dobra}. G-Eval to najbogatszy dostępny sygnał jakościowy
i~zarazem najdroższy oraz najbardziej podatny na bias. Architektura docelowa
(agent-sędzia w~Azure AI Foundry) jest zaprojektowana; brakuje rubryki i~dowodu,
że jej wyniki są warte swojej ceny.

\blok{Cel pracy}
Zaprojektować rubrykę G-Eval odwzorowaną na komponenty istniejącego silnika
jakości, zaimplementować ją jako analizator dostępny wyłącznie w~wariancie
chmurowym, i~zmierzyć --- względem etykiet ludzkich --- czy jej ocena jest
trafniejsza niż obecny wskaźnik telemetryczny, oraz ile to kosztuje.

\blok{Zakres prac}
\begin{itemize}
  \item Rubryka G-Eval odwzorowana na sześć komponentów silnika v2
        (\emph{reliability}, \emph{acceptance}, \emph{friction}, \emph{latency},
        \emph{feedback}, \emph{efficiency}) --- tak, by wynik dało się zestawić
        z~istniejącym score'em, a~nie dołożyć jako kolejną, niepowiązaną liczbę.
  \item Implementacja \texttt{IInsightAnalyzer} w~trybie „cloud-only'': lokalnie
        raportuje status \emph{no data}, nigdy błąd.
  \item Zbiór referencyjny 50--100 sesji (Seeder: \emph{clean}, \emph{error-prone},
        \emph{frustrated} + rozmowy scenariuszowe).
  \item Etykietowanie ręczne przez co najmniej dwie osoby, z~raportem zgodności
        międzysędziowskiej (Cohen $\kappa$) --- to jest \emph{ground truth}.
  \item Pomiar kosztu tokenowego na sesję jako funkcji długości transkryptu,
        z~wyznaczeniem progu opłacalności.
  \item Raport biasu sędziego: długość odpowiedzi jako fałszywe proxy jakości,
        bias pozycyjny, wrażliwość na sformułowanie rubryki.
\end{itemize}

\blok{Metodyka}
G-Eval występuje jako \emph{trzeci sędzia} porównywany do etykiet ludzkich ---
nie do istniejącego wskaźnika CopilotScope. Porównywanie modelu do liczby, którą
ma zastąpić, mierzy zgodność dwóch narzędzi, a~nie trafność żadnego z~nich.

\blok{Kryterium akceptacji}
Korelacja rangowa (Spearman) G-Eval~$\leftrightarrow$~etykiety ludzkie wyższa niż
korelacja obecnego wskaźnika 0--100~$\leftrightarrow$~etykiety, przy
udokumentowanym koszcie na sesję i~opublikowanej rubryce wraz z~przykładami
uzasadnień --- sama liczba bez rubryki jest w~tym projekcie bezużyteczna.

\blok{Punkt wejścia \quad·\quad Dlaczego magisterska}
\kod{Quality/Insights.cs} (interfejs), \kod{src/CopilotScope.AgentForge/Agents/}
(gotowy klient Azure AI Foundry), \kod{tools/CopilotScope.Seeder}. \\
Sednem pracy jest zaprojektowanie eksperymentu z~anotatorami i~analiza biasu,
a~nie implementacja --- ta jest tu najmniejszą częścią.

% ============================================================
\temat{2}{Uczenie interpretowalnych rubryk satysfakcji (SPUR) z~etykietowanych
sesji i~badanie ich transferu między asystentami}
{Praca magisterska}
{Algorytm \#2 --- SPUR (arXiv 2403.12388) \quad·\quad stan: nie zaimplementowany \quad·\quad wymaga zbioru etykietowanego}

\begin{streszczenie}
SPUR uczy model \emph{interpretowalnych} rubryk zadowolenia i~niezadowolenia
(SAT/DSAT) z~oznaczonych rozmów, zamiast oceniać według rubryki napisanej
z~góry, jak G-Eval. To stan sztuki w~szacowaniu satysfakcji użytkownika, badany
m.in.\ na Bing Copilot --- ale wymaga własnego zbioru etykiet, którego
CopilotScope nie ma. Praca bada, ile tych etykiet naprawdę trzeba, czy da się je
częściowo zastąpić słabym nadzorem, i~czy wyuczona rubryka przenosi się między
asystentami.
\end{streszczenie}

\blok{Kontekst i motywacja}
Temat \#1 (G-Eval) i~ten dzielą tę samą wąską gardło: brak etykiet. Różnią się
tym, co z~nimi robią --- G-Eval używa ich \emph{do walidacji} statycznej rubryki,
SPUR \emph{do jej wyuczenia}. Warto połączyć oba tematy w~jedną kampanię
etykietowania zamiast prowadzić dwie. Dodatkowo CopilotScope ma już
heurystyczny \texttt{FrustrationAnalyzer}, który --- mimo że zaszumiony ---
może posłużyć jako słaby etykietator wstępny.

\blok{Cel pracy}
Wyuczyć rubryki SAT/DSAT na sesjach CopilotScope, zmierzyć krzywą uczenia
względem liczby etykiet i~sprawdzić, czy rubryka wyuczona na jednym emiterze
telemetrii zachowuje trafność na innym.

\blok{Zakres prac}
\begin{itemize}
  \item Pilotaż etykietowania $\approx$200 sesji przez zespół, który sam używa
        CopilotScope na co dzień (dogfooding daje najwiarygodniejsze etykiety,
        bo etykietujący rozumie kontekst pracy).
  \item Wariant \emph{weak supervision}: \texttt{FrustrationAnalyzer} (algorytm \#9)
        jako wstępny, zaszumiony etykietator --- z~pomiarem, ile ręcznej pracy
        realnie oszczędza i~jaki bias wnosi.
  \item Krzywa uczenia: trafność SPUR w~funkcji liczby etykiet, z~wyznaczeniem
        punktu, od którego SPUR przebija statyczną rubrykę G-Eval na tym samym
        zbiorze walidacyjnym.
  \item Test transferu: rubryka wyuczona na sesjach VS~Code, oceniana na sesjach
        Copilot CLI i~Claude Code (trzy różne dialekty telemetrii, ten sam
        kolektor) oraz na innym repozytorium.
  \item Porównanie kosztu inferencji SPUR vs G-Eval --- SPUR nie wozi pełnej
        rubryki w~każdym promptcie, więc powinien być tańszy; trzeba to pokazać
        liczbami, a~nie założyć.
\end{itemize}

\blok{Metodyka}
Wspólny zbiór walidacyjny z~tematem \#1, żeby porównanie SPUR vs G-Eval było
uczciwe. Podział train/valid po \emph{sesjach}, nie po turach --- tury z~jednej
sesji są skorelowane i~wyciekłyby między zbiorami.

\blok{Kryterium akceptacji}
SPUR osiąga trafność nie gorszą niż G-Eval przy istotnie niższym koszcie na
inferencję, a~praca dostarcza udokumentowaną krzywą uczenia z~jasną odpowiedzią:
\emph{ile etykiet trzeba, zanim to się zaczyna opłacać}.

\blok{Ryzyko}
Rubryka wyuczona na jednym zespole i~jednym repozytorium może nie transferować
się nigdzie indziej. Test cross-repo jest tu obowiązkowy, a~nie opcjonalny ---
bez niego praca dowodzi tylko, że model nauczył się jednego zespołu.

\blok{Punkt wejścia}
\kod{Quality/FrustrationAnalyzer.cs} (słaby etykietator),
\kod{src/CopilotScope.AgentForge/}, \kod{research/07_frustration_analyzer.ipynb}.

% ============================================================
\temat{3}{Wykrywanie sesji opartych o retrieval w~telemetrii OTel GenAI jako
warunek wstępny metryk RAGAS}
{Praca inżynierska}
{Algorytm \#3 --- metryki komponentowe RAG (RAGAS) \quad·\quad stan: nie zaimplementowany \quad·\quad praca przygotowawcza}

\begin{streszczenie}
Metryki RAGAS --- \emph{faithfulness}, \emph{answer relevance},
\emph{context precision} --- mają sens wyłącznie wtedy, gdy sesja rzeczywiście
wykonuje retrieval: przeszukuje dokumentację repozytorium albo bazę wiedzy przed
odpowiedzią. Obecna telemetria CopilotScope w~ogóle nie odróżnia sesji
z~retrievalem od czysto generatywnej. Praca odpowiada na pytanie, które trzeba
zadać \emph{przed} napisaniem pierwszej linijki analizatora RAGAS: ile sesji
w~ogóle się do niego kwalifikuje?
\end{streszczenie}

\blok{Kontekst i motywacja}
RAGAS ma najmocniejsze zaplecze akademickie z~całej dziesiątki i~właśnie dlatego
jest najłatwiejszy do przecenienia. Dobre zaplecze teoretyczne nie oznacza
dopasowania do tego, co faktycznie robią sesje Copilota. Jeżeli okaże się, że
kwalifikuje się 5\% sesji, to jest to wynik --- oszczędza projektowi
zbudowania analizatora, który nigdy nie ma danych.

\blok{Cel pracy}
Ustalić, czy istniejące spany \texttt{execute\_tool} niosą wystarczający sygnał,
by rozpoznać retrieval bez nowej instrumentacji po stronie klienta, zmierzyć
odsetek kwalifikujących się sesji i~wydać rekomendację \emph{idź / nie idź}
popartą danymi.

\blok{Zakres prac}
\begin{itemize}
  \item Klasyfikator wywołań narzędzi na podstawie atrybutu
        \texttt{gen\_ai.tool.name}: retrieval (\texttt{search}, \texttt{grep},
        \texttt{read\_file}, \texttt{fetch}) kontra działanie
        (\texttt{run\_tests}, \texttt{execute\_command}, \texttt{edit\_file}) ---
        z~jawną listą przypadków niejednoznacznych.
  \item Pomiar: odsetek sesji z~co najmniej jednym wywołaniem sklasyfikowanym
        jako retrieval, osobno dla zbioru z~Seedera i~dla realnych sesji pracy
        nad repozytorium.
  \item Propozycja minimalnego atrybutu, który usunąłby heurystykę --- np.
        \kod{gen_ai.tool.category=retrieval} --- sformułowana jako
        \textbf{rozszerzenie konwencji semantycznych OTel GenAI}, a~nie prywatny
        atrybut CopilotScope. Prywatny atrybut rozwiązuje problem w~jednym
        narzędziu; propozycja do konwencji rozwiązuje go dla wszystkich.
  \item Prototyp jednej metryki (\emph{context precision}) na zakwalifikowanym
        podzbiorze --- dowód wykonalności, nie produkcyjny analizator.
  \item Raport \emph{idź / nie idź} z~rekomendacją priorytetu.
\end{itemize}

\blok{Metodyka}
Audyt danych przed implementacją. Klasyfikacja heurystyczna po nazwie narzędzia
walidowana ręcznie na próbce 100 wywołań (z~podaniem trafności heurystyki ---
jeśli sama heurystyka myli się w~30\% przypadków, to wynik audytu też jest
obarczony tym błędem i~trzeba to napisać).

\blok{Kryterium akceptacji}
Liczbowa odpowiedź na pytanie „jaki odsetek sesji kwalifikuje się do oceny
RAGAS'' wraz z~przedziałem ufności, plus jednoznaczna rekomendacja i~gotowa
propozycja atrybutu do zgłoszenia w~specyfikacji OTel GenAI.

\blok{Punkt wejścia}
\kod{src/CopilotScope.Collector/Otlp/} (dekoder spanów),
\kod{Domain/CopilotSession.cs}, \kod{tools/CopilotScope.Seeder}.

\blok{Jak podnieść do poziomu magisterskiego}
Zastąpić heurystykę po nazwie narzędzia klasyfikatorem uczonym na
zaetykietowanych wywołaniach i~wykazać przewagę --- wraz z~pełną implementacją
RAGAS na zakwalifikowanym podzbiorze.

% ============================================================
\temat{4}{Kalibracja modelu przeżywalności edycji (Edit Survival) względem
faktycznego stanu repozytorium git}
{Praca inżynierska}
{Algorytm \#4 --- Edit Survival Analysis \quad·\quad stan: zaimplementowany w~pełni \quad·\quad \texttt{EditSurvivalAnalyzer}}

\begin{streszczenie}
\texttt{EditSurvivalAnalyzer} łączy dwa sygnały --- pokrycie czterogramowe
i~brak cofnięcia edycji po 30/300~s --- wagami 0{,}4 i~0{,}6. Obie liczby
zostały przeniesione z~telemetrii GitHuba i~nigdy nie zostały zwalidowane na
danych CopilotScope. Praca konfrontuje wynik algorytmu z~faktem twardym
i~sprawdzalnym: czy wygenerowany kod nadal jest w~repozytorium tydzień później.
\end{streszczenie}

\blok{Kontekst i motywacja}
To najdojrzalszy z~dziesięciu algorytmów --- ma pełną implementację, notatnik
odtwarzający wszystkie wzory (\texttt{research/03\_edit\_survival.ipynb})
i~rozdział w~artykule wraz z~analizą wrażliwości wag. Czego nie ma, to
niezależnego punktu odniesienia. „Kod przeżył 300~sekund'' jest proxy dla
„kod był użyteczny''; git wie, czy przeżył naprawdę.

\blok{Cel pracy}
Zbudować narzędzie łączące edycje z~sesji z~historią git i~użyć go do
kalibracji wag oraz okien czasowych algorytmu.

\blok{Zakres prac}
\begin{itemize}
  \item Narzędzie mapujące edycje z~sesji na fragmenty kodu w~historii git
        (\texttt{git log}/\texttt{git blame} na plikach dotkniętych w~sesji).
  \item Zbiór 20--30 sesji z~twardym \emph{ground truth}: czy kod przetrwał
        w~repozytorium po 7 i~po 30 dniach.
  \item Test wrażliwości podziału wag 0{,}4/0{,}6 na siatce wartości --- artykuł
        (§\textit{Weight Sensitivity}) dostarcza gotowego aparatu, praca dokłada
        do niego rzeczywiste wyniki zamiast założeń.
  \item Weryfikacja okien 30/300~s osobno dla sesji \emph{inline} (VS~Code)
        i~\emph{agentowych} (Claude Code, Copilot CLI) --- rytm edycji w~długiej
        turze agentowej jest inny niż przy pojedynczej sugestii.
  \item Analiza przypadku, którego telemetria dziś nie rozróżnia: cofnięcie
        \emph{bo kod był zły} kontra cofnięcie \emph{bo zmienił się cel zadania}.
        Nawet negatywny wynik („nie da się rozróżnić bez nowego sygnału'') jest
        wartościowy, jeśli jest udokumentowany.
  \item Rekomendacja: utrzymać wagi / przestroić / dodać sygnał.
\end{itemize}

\blok{Kryterium akceptacji}
Udokumentowana korelacja między \emph{survival score} a~przeżyciem kodu
w~historii git, z~rekomendacją wag popartą liczbami --- oraz uczciwie
zaraportowanym rozmiarem próbki, bo 25 sesji to nie jest dowód ogólny.

\blok{Ryzyko}
Merytorycznie niskie --- algorytm jest dobrze opisany. Główne ryzyko to
nadinterpretacja małej próbki; praca powinna raportować przedziały ufności,
a~nie same punktowe współczynniki korelacji.

\blok{Punkt wejścia}
\kod{Quality/Analyzers.cs} (\kod{EditSurvivalAnalyzer}),
\kod{research/03_edit_survival.ipynb}, artykuł §\textit{Edit Survival Analysis}.

\blok{Jak podnieść do poziomu magisterskiego}
Zamienić prostą korelację na właściwy model przeżycia z~cenzurowaniem
(Kaplan--Meier, regresja Coxa) --- edycje w~sesjach trwających krócej niż okno
obserwacji są klasycznym przypadkiem danych cenzurowanych prawostronnie.

% ============================================================
\temat{5}{Czy przepustowość i~przeżywalność mierzą to samo? Analiza korelacyjna
metryki produktywności i~jej odporności na prawo Goodharta}
{Praca inżynierska}
{Algorytm \#5 --- Acceptance-weighted throughput \quad·\quad stan: zaimplementowany w~pełni \quad·\quad \texttt{ThroughputAnalyzer}}

\begin{streszczenie}
Przepustowość liczy zaakceptowane linie kodu na turę, dyskontowane
odrzuceniami --- prosty, celowo prosty proxy produktywności. Praca sprawdza,
czy ta prostota szkodzi: czy wysoka przepustowość idzie w~parze z~wysoką
przeżywalnością kodu (algorytm \#4), czy to dwie niezależne osie. Drugą częścią
pracy jest coś, czego zwykle się nie pisze, a~tutaj jest obowiązkowe: katalog
sposobów, na jakie tę metrykę da się zgamować.
\end{streszczenie}

\blok{Kontekst i motywacja}
Przepustowość nagradza \emph{dużo} zaakceptowanego kodu i~nie mówi nic o~tym,
czy \emph{mniej} kodu --- bo trafniejszego --- nie byłoby lepszym wynikiem.
Dokumentacja projektu nazywa ją wprost najbardziej podatną na prawo Goodharta
metryką z~całej dziesiątki („gdy miara staje się celem, przestaje być dobrą
miarą''). Wchodzi dziś do składanego wskaźnika z~wagą 0{,}20 w~komponencie
\emph{acceptance} --- warto wiedzieć, czy zasłużenie.

\blok{Cel pracy}
Ustalić ilościowo relację między przepustowością a~przeżywalnością, zaproponować
wariant metryki znormalizowany po złożoności zadania i~zbudować katalog
scenariuszy gamingu wraz z~propozycją ich detekcji.

\blok{Zakres prac}
\begin{itemize}
  \item Wykres rozrzutu przepustowość $\times$ przeżywalność na wspólnym zbiorze
        sesji; współczynniki Pearsona i~Spearmana wraz z~przedziałami ufności.
  \item Analiza podzbiorów: czy korelacja zmienia znak dla pewnych klas sesji
        (np.\ refaktoryzacja kontra pisanie nowego kodu). Ujemna korelacja
        w~jakimkolwiek podzbiorze jest mocnym argumentem, żeby przepustowość
        nigdy nie była samodzielnym wskaźnikiem sukcesu.
  \item Wariant „przepustowość na jednostkę złożoności zadania'' --- normalizacja
        po liczbie dotkniętych plików zamiast czystych linii kodu; wraz
        z~odpowiedzią, czy telemetria w~ogóle niesie sygnał złożoności.
  \item Wpływ stylu formatowania (automatyczny formatter kontra ręczne
        formatowanie) na porównywalność liczby linii między repozytoriami.
  \item Rozdział \textit{Jak tę metrykę zgamować}: co najmniej trzy
        udokumentowane scenariusze (np.\ akceptowanie wygenerowanego kodu
        hurtem, generowanie rozwlekłego kodu, sztuczne dzielenie tur) plus
        propozycja sygnału wykrywającego każdy z~nich.
  \item Rekomendacja dla wagi komponentu \emph{acceptance} w~silniku v2.
\end{itemize}

\blok{Kryterium akceptacji}
Ilościowa odpowiedź na pytanie „przepustowość i~przeżywalność mierzą to samo czy
coś innego'', rekomendacja wagi i~rozdział o~gamingu --- praca bez tego ostatniego
nie jest kompletna, choćby analiza statystyczna była bez zarzutu.

\blok{Punkt wejścia}
\kod{Quality/Analyzers.cs} (\kod{ThroughputAnalyzer}),
\kod{research/04_throughput.ipynb}, \kod{Quality/QualityEngine.cs} (wagi).

% ============================================================
\temat{6}{Kalibracja modelu tarcia i~napraw na poziomie tury (TFRA) --- badanie
anotacyjne i~uczenie wag kar z~danych}
{Praca magisterska}
{Algorytm \#6 --- Turn-level Friction \& Repair Analysis \quad·\quad stan: zaimplementowany w~pełni \quad·\quad \texttt{SegmentAnalyzer}}

\begin{streszczenie}
TFRA dzieli sesję na tury i~ocenia każdą z~osobna: karze błędy modelu
($-0{,}35\times$), błędy narzędzi ($-0{,}15\times$) i~pętle naprawcze
($-0{,}10$), a~latencję ocenia względem mediany \emph{tej samej} sesji, więc nie
karze wolnych modeli --- tylko tury odstające. To najlepiej zaprojektowany
algorytm z~dziesiątki, ale trzy wagi kar zostały dobrane ręcznie. Praca sprawdza,
czy odtwarzają one intuicję inżyniera czytającego transkrypt, i~czy da się je
wyuczyć zamiast ustawiać.
\end{streszczenie}

\blok{Kontekst i motywacja}
TFRA zasila komponent \emph{friction} (waga 0{,}20) w~składanym wskaźniku oraz
panel \textit{Turn analysis}, który wskazuje najlepszą i~najgorszą turę sesji
\emph{wraz z~powodami}. Jest interpretowalny z~konstrukcji, co czyni go dobrym
kandydatem do kalibracji na ludzkich ocenach: da się zestawić „algorytm ukarał
tę turę za X'' z~„człowiek uznał tę turę za złą, bo Y''.

\blok{Cel pracy}
Zmierzyć zgodność automatycznego wyniku TFRA z~ręczną anotacją tur i~ustalić,
czy wagi kar wyuczone z~danych biją wagi dobrane ręcznie.

\blok{Zakres prac}
\begin{itemize}
  \item Protokół anotacji: definicja etykiety („tura z~tarciem: tak/nie''
        plus powód), instrukcja dla anotatorów, procedura rozstrzygania sporów.
  \item Anotacja co najmniej 50 tur z~rozmów scenariuszowych dostępnych
        w~Seederze, przez co najmniej dwóch anotatorów, z~raportem zgodności
        ($\kappa$). Niska zgodność \emph{między ludźmi} jest wynikiem samym
        w~sobie: wyznacza sufit, powyżej którego żaden algorytm nie ma jak wejść.
  \item Porównanie TFRA z~anotacją: trafność, $\kappa$, oraz jakościowa analiza
        wszystkich przypadków rozbieżnych.
  \item Kalibracja progu „pętli naprawczej'': ile powtórzeń nieudanych wywołań
        odróżnia realne zapętlenie agenta od normalnego cyklu debugowania
        (spróbuj $\rightarrow$ nie działa $\rightarrow$ popraw), który wcale nie
        świadczy o~złej sesji.
  \item Zbadanie nieliniowości względem długości sesji: czy jedna zła tura
        w~sesji czterdziestoturowej powinna ważyć tyle samo, co w~trzyturowej.
  \item Wariant z~wagami wyuczonymi (optymalizacja pod zgodność z~anotacją)
        zestawiony z~obecnymi wagami ręcznymi --- wraz z~oceną, czy zysk
        z~uczenia jest wart utraty prostoty i~przenośności.
\end{itemize}

\blok{Kryterium akceptacji}
Zgodność (trafność albo $\kappa$) wariantu rekomendowanego z~ręczną anotacją
powyżej progu 0{,}7, wraz z~pełną listą przypadków rozbieżnych i~ich analizą.
Jeżeli wagi ręczne wypadną nie gorzej niż wyuczone --- to też jest wynik,
i~to dobry: broni prostoty obecnego rozwiązania.

\blok{Punkt wejścia}
\kod{Quality/SegmentAnalyzer.cs}, \kod{research/02_tfra_segment_analyzer.ipynb},
artykuł §\textit{Turn-level Friction \& Repair Analysis}.

\blok{Dlaczego magisterska}
Badanie anotacyjne z~kontrolą zgodności międzysędziowskiej plus uczenie
parametrów z~danych --- techniczna część implementacji jest tu najmniej
wymagająca.

% ============================================================
\temat{7}{Rekalibracja krzywej użyteczności opóźnienia dla sesji agentowych ---
czy progi 2~s i~8~s obowiązują programistów?}
{Praca inżynierska}
{Algorytm \#7 --- Latency-utility model \quad·\quad stan: zaimplementowany w~pełni \quad·\quad \texttt{LatencyUtilityAnalyzer}}

\begin{streszczenie}
Krzywa odwzorowująca czas do pierwszego tokenu na użyteczność opiera się na
progach z~ogólnych badań interakcji człowiek--komputer: 2~sekundy jako granica
utrzymania uwagi, 8~sekund jako granica porzucenia. Te progi nie pochodzą od
programistów czekających na sugestię kodu, a~już na pewno nie od użytkowników
agentów, którzy z~założenia akceptują wielosekundowe „myślenie'' narzędzia.
Praca konfrontuje krzywą z~niezależnym sygnałem: wykrytą frustracją użytkownika.
\end{streszczenie}

\blok{Kontekst i motywacja}
Komponent \emph{latency} ma wagę 0{,}15 w~składanym wskaźniku, a~pełny model
użyteczności działa jako osobny insight z~podziałem na kubełki ryzyka. Jeśli
progi są przeniesione z~niewłaściwej domeny, cały komponent systematycznie
przekłamuje --- w~jedną albo w~drugą stronę, w~zależności od tego, jakiego
asystenta używa zespół.

\blok{Cel pracy}
Ustalić empirycznie, czy progi uwagi i~porzucenia różnią się między sesjami
\emph{inline} a~\emph{agentowymi}, i~czy krzywa powinna być bezwzględna, czy
względna do mediany sesji.

\blok{Zakres prac}
\begin{itemize}
  \item Rozdzielenie sesji na inline (VS~Code) i~agentowe (Claude Code,
        Copilot CLI) --- klasyfikacja \texttt{SessionKind} istnieje w~kolektorze
        i~wymaga tylko rozszerzenia o~to rozróżnienie.
  \item Analiza krzyżowa z~algorytmem \#9: czy \texttt{FrustrationAnalyzer}
        wykrywa istotnie więcej sygnałów frustracji w~wiadomościach
        \emph{bezpośrednio po} wolnych turach niż w~losowo wybranych miejscach
        sesji. To jedyny sposób, żeby zewnętrznie zwalidować próg
        \emph{porzucenia} bez badania z~użytkownikami.
  \item Test efektu kontrastu: czy odbiór opóźnienia zależy od tego, jak szybka
        była poprzednia odpowiedź w~tej samej sesji --- czyli czy krzywa
        użyteczności powinna być względna do mediany sesji, tak jak robi to już
        TFRA przy ocenie latencji tury.
  \item Kalibracja progów osobno per rodzaj sesji i~per model --- artykuł
        wskazuje kalibrację per model jako jedną z~zaplanowanych ścieżek rozwoju.
  \item Oszacowanie, jaka część wariancji sygnału frustracji daje się wyjaśnić
        samą latencją, a~jaka pochodzi z~innych czynników.
\end{itemize}

\blok{Kryterium akceptacji}
Rekomendacja progów osobno dla sesji inline i~agentowych, poparta liczbowym
związkiem z~niezależnym sygnałem frustracji, oraz rozstrzygnięcie, czy krzywa ma
pozostać bezwzględna, czy przejść na wariant względny.

\blok{Ryzyko}
Progi z~badań interakcji człowiek--komputer są dobrze ugruntowane ogólnie.
Ryzykiem jest przestrojenie ich pod mały, niereprezentatywny zbiór sesji tego
projektu --- praca musi jasno powiedzieć, ile sesji stoi za rekomendacją.

\blok{Punkt wejścia}
\kod{Quality/Analyzers.cs} (\kod{LatencyUtilityAnalyzer}),
\kod{Quality/FrustrationAnalyzer.cs}, \kod{research/05_latency_utility.ipynb},
\kod{Domain/SessionKind.cs}.

% ============================================================
\temat{8}{Wrażliwość raportowanych kosztów i~oszczędności cache na nieaktualny
cennik --- analiza i~mechanizm kontroli świeżości}
{Praca inżynierska}
{Algorytm \#8 --- ekonomia tokenów i~cache \quad·\quad stan: zaimplementowany w~pełni \quad·\quad \texttt{TokenEconomicsAnalyzer}}

\begin{streszczenie}
Koszt sesji liczony jest z~cennika trzymanego w~konfiguracji
(\texttt{CopilotScope:Pricing}). Matematyka jest poprawna i~opisana w~artykule,
ale cała metryka opiera się na założeniu, którego nikt nie pilnuje: że cennik
jest aktualny. Dostawcy modeli zmieniają ceny częściej, niż ktokolwiek
aktualizuje plik konfiguracyjny. Praca mierzy, jak bardzo nieaktualny cennik
zniekształca wnioski, i~buduje mechanizm, który o~tym ostrzega.
\end{streszczenie}

\blok{Kontekst i motywacja}
To temat operacyjny, nie algorytmiczny --- i~właśnie dlatego łatwo go odkładać
w~nieskończoność, podczas gdy jego brak po cichu psuje jedną z~sześciu w~pełni
działających metryk. Metryka, która myli się o~40\% i~nie mówi o~tym ani słowa,
jest gorsza od metryki, której nie ma, bo ludzie podejmują na jej podstawie
decyzje.

\blok{Cel pracy}
Skwantyfikować wrażliwość raportowanego kosztu i~oszczędności cache na wiek
cennika, zaimplementować ostrzeżenie o~nieaktualności i~rozstrzygnąć dwie
otwarte kwestie definicyjne.

\blok{Zakres prac}
\begin{itemize}
  \item Przegląd historyczny cenników dwóch--trzech dostawców modeli
        z~ostatnich 6--12 miesięcy: jak często i~o~ile zmieniały się ceny wejścia,
        wyjścia i~odczytu z~cache.
  \item Symulacja: o~ile procent przesuwa się raportowany koszt sesji
        i~raportowana oszczędność z~cache, gdy cennik jest nieaktualny
        o~3 i~o~6 miesięcy. Wynik w~formie tabeli wrażliwości.
  \item Mechanizm „cennik nie był aktualizowany od $X$ dni'': ostrzeżenie
        widoczne w~raporcie insightu, nie w~logach --- log przeczyta programista,
        a~metrykę czyta menedżer.
  \item Zbadanie porównywalności między emiterami: czy oszczędności z~cache
        raportowane przez VS~Code i~przez Claude Code są w~ogóle porównywalne,
        czy różne strategie cache'owania po stronie dostawców czynią porównanie
        międzyemiterowe mylącym.
  \item Rozstrzygnięcie definicyjne: czy „koszt na zaakceptowaną edycję''
        powinien obejmować koszt \emph{odrzuconych} prób w~tej samej turze
        (koszt rzeczywisty), czy tylko koszt tury udanej (koszt krańcowy).
        Obie definicje są obronne --- praca ma wybrać jedną i~uzasadnić wybór.
\end{itemize}

\blok{Kryterium akceptacji}
Tabela wrażliwości pokazująca liczbowo skalę zniekształcenia, działające
ostrzeżenie o~świeżości cennika w~raporcie insightu oraz jednoznaczna,
uzasadniona rekomendacja definicji kosztu na edycję.

\blok{Punkt wejścia}
\kod{Quality/Analyzers.cs} (\kod{TokenEconomicsAnalyzer},
\kod{PricingOptions}), \kod{appsettings.json} (sekcja
\kod{CopilotScope:Pricing}), \kod{research/06_token_economics.ipynb}.

\blok{Jak podnieść do poziomu magisterskiego}
Rozszerzyć o~model kosztu alternatywnego: przy jakim stosunku ceny do jakości
opłaca się przełączyć sesję na tańszy model --- z~wykorzystaniem wyników
z~tematów \#4 i~\#5 jako miary jakości.

% ============================================================
\temat{9}{Dwujęzyczna (PL/EN) ewaluacja leksykalnego klasyfikatora frustracji
i~warunki jego awansu do składanego wskaźnika jakości}
{Praca magisterska}
{Algorytm \#9 --- klasyfikacja frustracji \quad·\quad stan: uproszczony, \emph{report-only} \quad·\quad \texttt{FrustrationAnalyzer}}

\begin{streszczenie}
\texttt{FrustrationAnalyzer} łączy leksykon polsko-angielski, miarę podobieństwa
Jaccarda do wykrywania przeformułowań i~sygnały typograficzne. Jest świadomie
zaszumiony i~dlatego świadomie trzymany \emph{poza} składanym wskaźnikiem
jakości --- raportuje, ale nie ocenia. Praca buduje zbiór walidacyjny, mierzy
precyzję i~czułość osobno dla polskiego i~angielskiego, i~ustala twardy próg,
po przekroczeniu którego sygnał wolno awansować do wskaźnika.
\end{streszczenie}

\blok{Kontekst i motywacja}
Dokumentacja projektu wymienia znane słabości tej heurystyki --- fałszywe
trafienia przy negacji („\emph{no worries}'', „\emph{nice, that works}''),
ślepotę na sarkazm, bias językowy --- ale nikt ich nie zmierzył. Ścieżka awansu
do składanego wskaźnika prowadzi przez walidację, nie przez dopisywanie kolejnych
reguł. Pokusa „wygląda wystarczająco dobrze, wrzućmy do score'u'' jest dokładnie
tym, przed czym ostrzega dokumentacja.

\blok{Cel pracy}
Zmierzyć rzeczywistą precyzję i~czułość obecnej heurystyki w~obu językach,
skatalogować jej systematyczne błędy i~sprawdzić, czy lekki klasyfikator uczony
przebija ją na tym samym zbiorze walidacyjnym.

\blok{Zakres prac}
\begin{itemize}
  \item Korpus walidacyjny: 100+ wiadomości z~sesji persony \emph{frustrated}
        oraz 100 z~persony \emph{clean} jako kontrola negatywna, uzupełnione
        wiadomościami z~realnych sesji.
  \item Etykietowanie przez co najmniej dwie osoby, z~raportem zgodności
        ($\kappa$). Frustracja jest zjawiskiem subiektywnym --- bez miary
        zgodności między ludźmi nie da się uczciwie ocenić algorytmu.
  \item Macierz pomyłek \textbf{osobno dla języka polskiego i~angielskiego}.
        Bias językowy jest podejrzewany od początku, ale niezmierzony;
        leksykon polski jest mniejszy i~może mieć zupełnie inną charakterystykę.
  \item Katalog błędów systematycznych z~przykładami: negacja, sarkazm, krótkie
        entuzjastyczne zdania, wtrącenia techniczne. Udokumentowane ograniczenie
        jest wartościowym wynikiem --- ukryte jest wadą.
  \item Lekki klasyfikator uczony (regresja logistyczna na reprezentacjach
        wektorowych albo mały model językowy) jako konkurent heurystyki,
        z~porównaniem trafności i~kosztu.
  \item Propozycja progu awansu w~formie zdania rozstrzygalnego, np.\
        „promocja do składanego wskaźnika wymaga precyzji $\geq X$ przy czułości
        $\geq Y$, mierzonych osobno w~każdym z~języków''.
\end{itemize}

\blok{Kryterium akceptacji}
Udokumentowany baseline (precyzja/czułość/F1) obecnej heurystyki dla obu języków
oraz jawny, liczbowy próg awansu --- zamiast subiektywnej oceny „wygląda nieźle''.

\blok{Wymiar etyczny}
To jedyny z~dziesięciu algorytmów, który dotyka emocji człowieka wprost. Praca
musi utrzymać zasadę \emph{report-only} dopóki walidacja nie jest jednoznaczna
i~zawierać rozdział o~ryzyku nadużycia --- wskaźnik frustracji użyty do oceny
\emph{ludzi}, a~nie sesji, jest dokładnie tym zastosowaniem, przed którym projekt
się broni brakiem tożsamości w~danych.

\blok{Punkt wejścia}
\kod{Quality/FrustrationAnalyzer.cs}, \kod{research/07_frustration_analyzer.ipynb},
artykuł §\textit{Lexical Frustration Classifier},
\kod{tools/CopilotScope.Seeder/Personas.cs}.

% ============================================================
\temat{10}{Detekcja domknięcia zadania na podstawie zewnętrznych sygnałów CI ---
prototyp end-to-end od sesji do dashboardu}
{Praca inżynierska}
{Algorytm \#10 --- task-completion detection \quad·\quad stan: nie zaimplementowany, częściowe hooki \quad·\quad wymaga sygnałów spoza telemetrii}

\begin{streszczenie}
To pytanie najbliższe temu, o~co naprawdę chodzi --- \emph{czy praca została
wykonana} --- i~zarazem najtrudniejsze do zautomatyzowania z~całej dziesiątki.
Telemetria asystenta sama z~siebie nie mówi, czy build przeszedł ani czy testy
są zielone. Praca buduje najprostszą działającą ścieżkę: krok w~systemie
ciągłej integracji wysyła wynik budowania i~testów do istniejącego API kolektora,
a~dashboard pokazuje sesję jako domkniętą albo urwaną.
\end{streszczenie}

\blok{Kontekst i motywacja}
Wszystkie pozostałe dziewięć algorytmów mierzą \emph{przebieg} sesji: ile było
błędów, jak szybko odpowiadał model, ile kodu przetrwało. Żaden nie odpowiada na
pytanie, czy zadanie zostało domknięte. Sygnału trzeba szukać poza telemetrią
asystenta --- w~kodzie wyjścia budowania, wyniku testów i~stanie repozytorium po
sesji. Kolektor ma już API przyjmujące dane; brakuje kroku po stronie CI
i~komponentu, który z~tego sygnału coś wyliczy.

\blok{Cel pracy}
Zbudować i~zwalidować kompletną ścieżkę od zakończenia sesji do widocznego
w~dashboardzie sygnału domknięcia, oraz zmierzyć, jak wiele sesji w~ogóle taki
sygnał otrzymuje.

\blok{Zakres prac}
\begin{itemize}
  \item Ustalenie minimalnego zestawu sygnałów wystarczającego, by odróżnić
        „zadanie domknięte'' od „sesja się urwała'': kod wyjścia budowania,
        wynik testów, czysty stan repozytorium po sesji.
  \item Rozszerzenie API kolektora o~endpoint przyjmujący sygnały domknięcia
        oraz analizator \texttt{IInsightAnalyzer}, który je interpretuje.
  \item Krok w~GitHub Actions wysyłający po \texttt{dotnet test} ładunek
        \kod{{sessionId, buildPassed, testsGreen}}. \textbf{Świadomie jeden
        system CI} --- ten, którego repozytorium już używa. Projektowanie
        uniwersalnego adaptera na wszystkie systemy CI jest tu główną pułapką
        i~pracą, która nigdy się nie kończy.
  \item Pomiar: jaki odsetek sesji kończy się bez żadnego zewnętrznego sygnału
        domknięcia --- i~rzetelna odpowiedź, czy „programista po prostu przestał
        pisać'' da się w~ogóle odróżnić od „zadanie skończone, tylko nikt tego
        nie zaraportował'' bez sędziego LLM (temat \#1).
  \item Walidacja przez dogfooding: 5--10 realnych sesji pracy nad tym
        repozytorium, przeprowadzonych od zera do zielonego CI.
  \item Prezentacja sygnału w~panelu \textit{Insights} wraz z~poprawną obsługą
        stanu \emph{no data} --- brak sygnału nie może obniżać oceny sesji.
\end{itemize}

\blok{Kryterium akceptacji}
Działający przykład end-to-end: sesja $\rightarrow$ pull request $\rightarrow$ CI
$\rightarrow$ sygnał domknięcia widoczny w~dashboardzie --- zademonstrowany na
realnych sesjach, zanim ktokolwiek zaprojektuje wariant chmurowy z~agentem-sędzią.

\blok{Ryzyko}
Największe ryzyko projektowe z~dziesięciu: łatwo zbudować rozwiązanie wymagające
integracji z~każdym systemem CI z~osobna, które zostaje kolejnym niedokończonym
adapterem. Zakres pracy celowo ogranicza się do jednego systemu; uogólnienie jest
tematem na później i~tylko wtedy, gdy wariant jednosystemowy okaże się użyteczny.

\blok{Punkt wejścia}
\kod{src/CopilotScope.Collector/Program.cs} (API ingest),
\kod{Quality/Insights.cs}, \kod{.github/workflows/ci.yml},
\kod{src/CopilotScope.Dashboard/Components/Pages/}.

% ============================================================
\clearpage
\section*{Dokumenty źródłowe i~dalsza lektura}
\addcontentsline{toc}{section}{Dokumenty źródłowe i dalsza lektura}

\begin{itemize}
  \item \texttt{research/articles/quality\_measurement\_framework.tex} ---
        formalny opis ośmiu zaimplementowanych algorytmów: definicje, wzory,
        własności, przykłady liczbowe krok po kroku oraz analiza wrażliwości wag.
        PDF jest dołączany do każdego wydania projektu.
  \item \texttt{docs/ANALYSIS.md} §8 --- przegląd dziesięciu podejść do oceny
        jakości sesji wraz z~zastrzeżeniami wobec każdego; §8a --- macierz
        implementacji; §8b --- referencja komponentów silnika v2.
  \item \texttt{README.md}, tabela \textit{Evaluation algorithms --- implementation
        status} --- ten sam zestaw dziesięciu algorytmów z~podziałem na warianty
        lokalny i~chmurowy. Ta sama numeracja obowiązuje w~niniejszym dokumencie.
  \item \texttt{research/RESEARCH\_PROPOSALS.md} --- wcześniejsza, krótsza wersja
        tych propozycji w~formie notatki roboczej; niniejszy dokument ją zastępuje
        jako materiał do przedstawienia studentom.
  \item \texttt{research/0X\_*.ipynb} --- notatniki Pythona odtwarzające każdy
        wzór i~przykład liczbowy z~artykułu; naturalny punkt startowy dla części
        eksperymentalnej większości tych prac.
  \item \texttt{docs/AGENTFORGE.md} --- opis eksperymentalnej usługi agentów
        person opartych o~sesje za zgodą; istotny kontekst dla tematów \#1 i~\#2,
        bo tam znajduje się działający klient Azure AI Foundry.
  \item Literatura zewnętrzna: Liu i~in., \textit{G-Eval: NLG Evaluation using
        GPT-4 with Better Human Alignment} (2023); \textit{SPUR --- Supervised
        Prompting for User satisfaction Rubrics}, arXiv:2403.12388;
        Es i~in., \textit{RAGAS: Automated Evaluation of Retrieval Augmented
        Generation} (2023); konwencje semantyczne OpenTelemetry GenAI.
\end{itemize}

\end{document}
